\documentclass{article}
\usepackage[T1]{fontenc}
\usepackage{lmodern}
\usepackage{graphicx}
\usepackage{amsmath,amssymb}
\usepackage[a4paper,margin=2cm]{geometry}
\usepackage[absolute,overlay]{textpos}
\setlength{\TPHorizModule}{1mm}
\setlength{\TPVertModule}{1mm}
\usepackage[indonesian]{babel}
\usepackage{hyperref}
\setlength{\emergencystretch}{2em}
% unit: Halaman 12

\begin{document}

% === Halaman 12 (llm) ===

\begin{itemize}
    \item Selanjutnya Alice menghitung kunci privat $d$ dengan kekongruenan:
    \[ ed \equiv 1 \pmod{\phi(n)} \longrightarrow d = \frac{1 + k \phi(n)}{e} \]
    $d$ adalah balikan $e$ dalam modulus $\phi(n)$ \\
    $d$ dapat dihitung dengan dengan rumus:
    \[ d = \frac{1 + k\phi(n)}{e} = \frac{1 + 3220k}{79} \]
    \vspace{5mm}
    Coba berbagai nilai-nilai $k = 1, 2, 3, \dots$, sampai mendapatkan $d$ bilangan bulat:
    \begin{align*}
        k = 1 & \rightarrow d = (1 + 3220 \times 1)/79 = 3221/79 = 40,77 & \text{(bukan bilangan bulat)}
        k = 2 & \rightarrow d = (1 + 3220 \times 2)/79 = 6441/79 = 81,53 & \text{(bukan bilangan bulat)}
        k = \dots & \text{dst}
    \end{align*}
    \textbf{k = 25 $\rightarrow$ d = (1 + 3220 $\times$ 25)/79 = 80501/79 = 1019 (bilangan bulat)}
    \newline diperoleh nilai $d$ yang bulat adalah 1019. Ini adalah kunci privat.
    \item Jadi, kunci privat Alice adalah $(d, n) = (1019, 3337)$. \textcolor{red}{Keep them secret!}
\end{itemize}

\end{document}
